\documentclass[11pt,a4paper]{article}

\usepackage[margin=2.3cm]{geometry}
\usepackage[T1]{fontenc}
\usepackage[utf8]{inputenc}
\usepackage{lmodern}
\usepackage{amsmath,amssymb,bm}
\usepackage{booktabs}
\usepackage{enumitem}
\usepackage{xcolor}
\usepackage{listings}
\usepackage{hyperref}

\hypersetup{
  colorlinks=true,
  linkcolor=blue,
  urlcolor=blue
}

\definecolor{codegray}{RGB}{245,245,245}
\definecolor{codeblue}{RGB}{0,70,140}

\lstset{
  language=Python,
  basicstyle=\ttfamily\small,
  backgroundcolor=\color{codegray},
  keywordstyle=\color{codeblue},
  frame=single,
  breaklines=true,
  showstringspaces=false
}

\title{\textbf{Practical Session: Locomotion Software for Legged Robots in Simulation}}
\author{}
\date{}

\begin{document}

\maketitle

\section{Objectives}

This practical session introduces the main components of a model-based locomotion software stack for a quadruped robot in simulation. The goal is to progressively build, tune, and test a controller that combines centroidal Model Predictive Control (MPC), gait scheduling, foothold generation, swing-leg control, stance-force control, inverse kinematics, state estimation, and terrain adaptation.

The experiments can be performed in MuJoCo using the provided robot model and simulation interface.

\section{Centroidal Model in CasADi}

The centroidal dynamics describe the evolution of the robot center of mass (CoM) and angular momentum under the action of contact forces. For a robot of mass \(m\), CoM position \(\bm{c}\), and contact forces \(\bm{f}_i\), the translational dynamics are:

\begin{align*}
m\ddot{\bm{c}} &= \sum_{i=1}^{n_c} \bm{f}_i + m\bm{g}.
\end{align*}

The angular momentum dynamics about the CoM are:

\begin{align*}
\dot{\bm{L}} &=
\sum_{i=1}^{n_c}
(\bm{p}_i-\bm{c}) \times \bm{f}_i,
\end{align*}

where \(\bm{p}_i\) is the position of the \(i\)-th contact point.

A possible state is:

\begin{align*}
\bm{x} =
\begin{bmatrix}
\bm{c} &
\dot{\bm{c}} &
\bm{\theta} &
\bm{\omega}
\end{bmatrix}^{T},
\end{align*}

while the MPC input contains the contact forces:

\begin{align*}
\bm{u} =
\begin{bmatrix}
\bm{f}_1 &
\bm{f}_2 &
\bm{f}_3 &
\bm{f}_4
\end{bmatrix}^{T}.
\end{align*}

\subsection*{Exercise}

\begin{enumerate}[label=\arabic*.]
    \item Implement the continuous centroidal dynamics in CasADi.
    \item Discretize the model using forward Euler or Runge--Kutta integration.
    \item Define a finite-horizon optimal control problem.
    \item Add unilateral-contact and friction constraints.
    \item Compare the centroidal prediction with the simulated CoM evolution in MuJoCo.
\end{enumerate}

For each active contact:

\begin{align*}
f_{z,i} &\geq 0, \\
|f_{x,i}| &\leq \mu f_{z,i}, \\
|f_{y,i}| &\leq \mu f_{z,i}.
\end{align*}



\section{Raibert-Style Foothold Generator}

At lift-off, the desired touchdown position can be computed from the nominal hip location, desired velocity, and velocity tracking error:

\begin{align*}
\bm{p}_{f,i}^{\mathrm{des}} =
\bm{p}_{\mathrm{hip},i}
+
\bm{p}_{i}^{\mathrm{nom}}
+
\frac{T_{\mathrm{st}}}{2}\bm{v}^{\mathrm{des}}
+
k_v
\left(
\bm{v} - \bm{v}^{\mathrm{des}}
\right),
\end{align*}

where \(T_{\mathrm{st}}\) is the stance duration and \(k_v\) is the velocity-feedback gain.


\subsection*{Exercises}

\begin{enumerate}[label=\arabic*.]
    \item Implement the nominal foothold generator.
    \item Command forward, lateral, and turning motion.
    \item Visualize the planned touchdown locations.
    \item Tune \(k_v\) and observe its influence on velocity recovery.
    \item Add reachable-workspace bounds around the hip.
\end{enumerate}

\section{Swing-Foot Trajectory}

The swing-foot trajectory must move the foot smoothly from lift-off to the desired touchdown point. It should provide sufficient clearance above the terrain and have low touchdown velocity.

A Cartesian feedback-linearization law is:

\begin{align*}
\ddot{\bm{p}}^{\mathrm{cmd}} =
\ddot{\bm{p}}^{\mathrm{ref}}
+
\bm{K}_p
\left(
\bm{p}^{\mathrm{ref}}-\bm{p}
\right)
+
\bm{K}_d
\left(
\dot{\bm{p}}^{\mathrm{ref}}-\dot{\bm{p}}
\right).
\end{align*}

Implement and compare:

\begin{enumerate}[label=\arabic*.]
    \item A Cartesian feedback-linearization controller;
    \item A cubic polynomial trajectory implemented numerically;
    \item A cubic spline generated with \texttt{scipy.interpolate}.
\end{enumerate}

For a cubic trajectory:

\begin{align*}
p(t) = a_0 + a_1t + a_2t^2 + a_3t^3.
\end{align*}

Choose the coefficients to satisfy initial and final position and velocity boundary conditions.

\subsection*{Exercises}

\begin{enumerate}[label=\arabic*.]
    \item Generate horizontal motion from lift-off to touchdown.
    \item Add a vertical swing-height profile.
    \item Compare position, velocity, and acceleration continuity.
    \item Change the swing height and test obstacle clearance.
    \item Verify that touchdown velocity is close to zero.
\end{enumerate}

\section{Stance Controller}

During stance, the desired contact force computed by the MPC is converted into joint torques through the foot Jacobian:

\begin{align*}
\bm{\tau}_i = \bm{J}_i^T \bm{f}_i^{\mathrm{des}}.
\end{align*}

The Jacobian must be consistent with the frame in which the contact force is expressed.

\subsection*{Exercises}

\begin{enumerate}[label=\arabic*.]
    \item Extract the foot Jacobian from the robot model.
    \item Apply a desired vertical force to one leg.
    \item Check the resulting joint torques.
    \item Compare desired and measured ground-reaction forces.
    \item Verify force direction and coordinate-frame conventions.
\end{enumerate}

\section{Inverse Kinematics}

Inverse kinematics maps desired swing-foot trajectories into desired joint positions or velocities.

\subsection{Numerical Inverse Kinematics}

A damped least-squares solution is:

\begin{align*}
\dot{\bm{q}} =
\bm{J}^{T}
\left(
\bm{J}\bm{J}^{T}+\lambda^2\bm{I}
\right)^{-1}
\dot{\bm{p}}^{\mathrm{des}}.
\end{align*}

The joint reference can then be obtained by numerical integration:

\begin{align*}
\bm{q}_{k+1}^{\mathrm{des}}
=
\bm{q}_k^{\mathrm{des}}
+
\dot{\bm{q}}^{\mathrm{des}} \Delta t.
\end{align*}

\subsection{QP-Based Inverse Kinematics}

A QP formulation is:

\begin{align*}
\min_{\dot{\bm{q}}}
\quad &
\left\|
\bm{J}\dot{\bm{q}}-\dot{\bm{p}}^{\mathrm{des}}
\right\|^2
+
\lambda
\left\|
\dot{\bm{q}}
\right\|^2,
\end{align*}

subject to joint limits and velocity bounds.

\subsection*{Exercises}

\begin{enumerate}[label=\arabic*.]
    \item Implement numerical IK for one leg.
    \item Add joint limits and damping near singularities.
    \item Implement or use a QP-based IK solver.
    \item Compare Cartesian tracking errors for the two methods.
    \item Use the resulting joint references during swing.
\end{enumerate}


\section{MPC Tuning}

The MPC should track desired CoM motion, body orientation, velocity, and contact forces. A typical cost function is:

\begin{align*}
	J =
	\sum_{k=0}^{N-1}
	\left[
	\|\bm{x}_k-\bm{x}^{\mathrm{ref}}_k\|_{\bm{Q}}^2
	+
	\|\bm{u}_k-\bm{u}^{\mathrm{ref}}_k\|_{\bm{R}}^2
	+
	\|\bm{u}_k-\bm{u}_{k-1}\|_{\bm{R}_{\Delta}}^2
	\right]
	+
	\|\bm{x}_{N}-\bm{x}^{\mathrm{ref}}_{N}\|_{\bm{Q}_f}^2.
\end{align*}

Important parameters are:

\begin{itemize}
	\item Prediction horizon \(N\);
	\item MPC sampling time;
	\item Contact schedule and selected gait;
	\item CoM position and velocity weights;
	\item Base orientation and angular velocity weights;
	\item Contact-force regularization and force-rate regularization;
	\item Friction coefficient and force limits.
\end{itemize}

\subsection*{Exercises}

\begin{enumerate}[label=\arabic*.]
	\item Start from a standing robot and regulate the CoM to a fixed reference.
	\item Command a constant forward velocity.
	\item Change the prediction horizon and evaluate stability, tracking, and computational time.
	\item Increase the velocity tracking weight and observe the resulting gait.
	\item Increase force regularization and observe the reduction in force variation.
	\item Compare the behaviour for different gait frequencies.
	\item Try  to reduce the MPC replanning frequency, what happens to the robot behaviour?
\end{enumerate}

\section{Gait Scheduler}

The gait scheduler generates a binary contact state for each foot:

\begin{align*}
	s_i(t) \in \{0,1\},
\end{align*}

where \(s_i=1\) denotes stance and \(s_i=0\) denotes swing.

Each gait is described by its cycle duration, duty factor, and phase offsets.

\begin{table}[h]
	\centering
	\begin{tabular}{lll}
		\toprule
		Gait & Swing pattern & Typical use \\
		\midrule
		Crawl & One foot swings at a time & Slow and statically stable walking \\
		Pace & Left and right lateral pairs alternate & Fast lateral-pair gait \\
		Trot & Diagonal pairs alternate & Dynamic walking and running \\
		Bound & Front and hind pairs alternate & Fast locomotion \\
		\bottomrule
	\end{tabular}
\end{table}

\subsection*{Exercises}

\begin{enumerate}[label=\arabic*.]
	\item Implement a periodic phase variable.
	\item Generate the contact schedule for crawl, pace, and trot.
	\item Plot the contact state of each leg over several gait periods.
	\item Ensure that the gait scheduler, MPC contacts, and swing-leg state are synchronized.
	\item Change duty factor and gait frequency during simulation.
\end{enumerate}

\section{Foothold Optimization and Push Recovery}

Foothold optimization modifies the nominal foothold generated by the Raibert rule to improve balance recovery and maintain feasible steps.

A simple objective is:

\begin{align*}
\min_{\bm{p}_f}
\quad &
\left\|
\bm{p}_f-\bm{p}_f^{\mathrm{nom}}
\right\|_{\bm{W}}^2
+
w_c
\left\|
\bm{p}_f-\bm{p}_{\mathrm{capture}}
\right\|^2,
\end{align*}

subject to feasible workspace and terrain constraints.

\subsection*{Experiment}

\begin{enumerate}[label=\arabic*.]
    \item Walk at constant velocity with foothold optimization disabled.
    \item Apply a push to the robot in MuJoCo (drag the robot with the mouse).
    \item Record CoM position, velocity error, and planned footholds.
    \item Enable foothold optimization.
    \item Repeat the same push and compare the recovery.
\end{enumerate}

The expected behaviour is a corrective foothold placed in the direction required to recover balance.

\section{Integral Action and Payload Adaptation}

Integral action can compensate for persistent tracking errors caused by unmodelled disturbances, such as an additional payload:

\begin{align*}
\bm{e}_I(t)
=
\int_{0}^{t}
\left(
\bm{v}^{\mathrm{des}}-\bm{v}
\right)
dt.
\end{align*}

The integral state can be added to the body-position or velocity controller.

\subsection*{Experiment}

\begin{enumerate}[label=\arabic*.]
    \item Tune the controller for the nominal robot.
    \item Add an additional 5kg mass to the robot model in MuJoCo.
    \item Test walking without integral action.
    \item Activate the integral term and compare velocity tracking.
\end{enumerate}

\section{State Estimation}

Replace the perfect simulated state with an estimated state. The estimator may combine:

\begin{itemize}
    \item IMU angular velocity and linear acceleration;
    \item Joint encoders;
    \item Foot kinematics;
    \item Contact detection;
    \item Optional ground-truth state only for debugging.
\end{itemize}

An Extended Kalman Filter can estimate base orientation, position, velocity, and IMU biases.

\subsection*{Exercises}

\begin{enumerate}[label=\arabic*.]
    \item Compare estimated and ground-truth base velocity.
    \item Add noise to IMU and joint measurements.
    \item Tune process noise and measurement noise.
    \item Disable contact updates and observe estimator drift.
    \item (optional) Study the effect of incorrect contact detection. Consider implementing Haptic touch-down.
\end{enumerate}

\section{Frontal Reflex}

A frontal reflex can be activated when the tracking error exceeds a threshold:

\begin{align*}
\left\|
\bm{e}_{\mathrm{track}}
\right\|
>
e_{\mathrm{th}}.
\end{align*}

Possible reflex actions are:

 

\subsection*{Exercises}

\begin{enumerate}[label=\arabic*.]
    \item Define a frontal tracking error based on pitch, CoM displacement, or velocity.
    \item Select an activation threshold.
    \item Move against an obstacle (e.g. a step) with the blind robot
    \item Verify that the reflex is inactive during nominal walking.
    \item Tune the threshold to avoid late or unnecessary activation.
\end{enumerate}

\section{Vision-Based Foothold Height Adjustment}

Use a terrain-height estimate to modify the vertical component of the planned foothold:

\begin{align*}
p_{f,z}^{\mathrm{des}} =
h_{\mathrm{terrain}}
\left(
p_{f,x}^{\mathrm{des}},
p_{f,y}^{\mathrm{des}}
\right)
+
h_{\mathrm{margin}}.
\end{align*}

The safety margin accounts for terrain-map uncertainty and contact-model imperfections.

\subsection*{Exercises}

\begin{enumerate}[label=\arabic*.]
    \item Test the controller on flat terrain.
    \item Add steps, ramps, and isolated obstacles.
    \item Query the terrain height at the planned foothold.
    \item Adjust touchdown height and swing-foot clearance.
    \item Compare missed contacts with and without height adjustment.
\end{enumerate}

\section{Rough-Terrain Tests and Terrain Estimation}

On uneven terrain, estimate a local terrain plane from the stance-foot positions. The terrain normal can then be used to generate a terrain-aligned desired base orientation.

This adaptation helps maintain a suitable leg configuration and improves contact consistency on slopes and uneven ground.

\begin{table}[h]
\centering
\begin{tabular}{p{0.32\textwidth}p{0.56\textwidth}}
\toprule
Configuration & Expected behaviour \\
\midrule
Fixed horizontal body reference &
The body remains aligned with the world frame; leg workspace and contact quality may degrade on slopes. \\[0.2cm]
Terrain-aligned body reference &
The desired body roll and pitch follow the estimated local terrain plane. \\[0.2cm]
Terrain-aware orientation and footholds &
The robot adapts both body pose and foot touchdown height, improving robustness on rough terrain. \\
\bottomrule
\end{tabular}
\end{table}

\subsection*{Final Experiment}

\begin{enumerate}[label=\arabic*.]
    \item Create a rough-terrain scenario in MuJoCo.
    \item Run the robot with terrain estimation disabled.
    \item Repeat with terrain estimation enabled.
    \item Compare base orientation, foot tracking error, contact forces, and velocity tracking.
    \item Discuss the role of terrain estimation in stability and locomotion robustness.
\end{enumerate}

%\section{Expected Final Demonstration}
%
%At the end of the practical session, the robot should be able to:
%
%\begin{itemize}
%    \item Walk using crawl, pace, and trot gaits;
%    \item Track a commanded body velocity using centroidal MPC;
%    \item Generate swing-foot trajectories and stance torques;
%    \item Use numerical or QP-based inverse kinematics;
%    \item Adapt footholds with Raibert feedback and foothold optimization;
%    \item Recover from pushes through corrective stepping;
%    \item Compensate partially for an added payload using integral action;
%    \item Walk on uneven terrain using terrain-aware foothold and base-orientation references.
%\end{itemize}

\end{document}
